\subsection{ソフトウェア工学プロセスの基礎}
\subsubsection{導入}
ソフトウェア工学プロセスとは、ソフトウェアを構築し、運用するためにソフトウェア技術者が実施する作業活動である。一般的な工学プロセスは、資源を消費しながら、入力を出力へ変換する相互に関連した活動の集合である。ソフトウェア工学も工学の一分野であるため、設計、製造、運用といった他の工学分野の考え方から多くを学んできた。

ただし、ソフトウェアでは、化学反応器のように同一の物理製品を多数製造するわけではない。むしろ、多数のソフトウェア単位を構築し、それらの流れを管理する点で製造の比喩が役立つ。さらに、ソフトウェア単位の実行は、ある種類のデータを別の種類のデータへ変換する操作として理解できる。

この章で使う「プロセス」という語は、プログラムが実行時に動く処理ではなく、技術者やチームが行う作業活動を意味する。プロセスは多くの知識領域と関係するが、とくにソフトウェア工学マネジメント、モデルと手法、品質、アーキテクチャ、テスト、測定、根本原因分析と深く関係する。

\begin{pointbox}{重要}どのプロジェクトにもそのまま適用できる理想プロセスは存在しない。プロセスは、プロジェクトと組織の文脈に合わせて選択し、適応し、測定しながら運用するものである。\end{pointbox}
\par\smallskip
\begin{center}
\centering
\IfFileExists{assets/generated_figures/ch10/fig-ch10-02.png}{\includegraphics[width=0.96\linewidth]{assets/generated_figures/ch10/fig-ch10-02.png}}{\fbox{\parbox[c][48mm][c]{0.9\linewidth}{\centering 画像生成待ち\\プロセスの入力・活動・出力の図}}}
\par\smallskip
\refstepcounter{figure}\label{fig:ch10-02}図 \thefigure: プロセスの入力・活動・出力の図
\end{center}
図 \thefigure は、プロセスの入力・活動・出力の図を視覚的に整理したものである。
\par\smallskip
\subsubsection{ソフトウェア工学プロセスの定義}
プロセスは、相互に関係し合う活動の集合であり、入力を出力へ変換する。活動は、まとまりのあるタスクの集合である。タスクは、一つ以上のプロセス成果を達成するために必要、推奨、または許容される行動である。

プロセスの説明には、通常、必要な入力、変換活動、生成される出力が含まれる。さらに、指示、規則、制約のような制御要素と、ツール、技術、人員、インフラのような支援要素も関係する。あるプロセスの出力は、別のプロセスの入力になることが多い。

ソフトウェアシステムは、ソフトウェアだけでなく、ハードウェア、人、手作業の手順を含むことがある。したがって、ソフトウェア工学プロセスを考えるときは、ソフトウェア部分だけでなく、その周辺の人と組織の動きも含めて捉える必要がある。

\paragraph{表3：プロセスを記述するときに見る要素}
\par\smallskip\noindent\refstepcounter{table}\label{tab:ch10-03}表 \thetable: 表3：プロセスを記述するときに見る要素における要素・意味・例\par\smallskip
表 \thetable は、表3：プロセスを記述するときに見る要素における要素・意味・例を比較・確認しやすい形で整理したものである。\par\smallskip
\begin{longtable}{>{\raggedright\arraybackslash}p{0.16\linewidth}>{\raggedright\arraybackslash}p{0.30\linewidth}>{\raggedright\arraybackslash}p{0.46\linewidth}}
\toprule
要素 & 意味 & 例\\
\midrule
\endfirsthead
\toprule
要素 & 意味 & 例\\
\midrule
\endhead
入力 & 活動の前提になる成果物や情報である。 & 要求、標準、既存コード、設計方針。\\
活動 & 入力を変換する作業である。 & 分析、設計、実装、レビュー、テスト。\\
出力 & 活動の結果として作られる成果物である。 & 設計書、コード、テスト結果、運用手順。\\
制御 & 活動を縛る方針や条件である。 & 契約、規制、品質基準、納期制約。\\
支援機構 & 活動を可能にする資源である。 & 人員、ツール、開発環境、組織知識。\\
\bottomrule
\end{longtable}
